\chapter
[\CO{[PSY]} \emph{Spezielle Patientengruppe:} Psychiatrie]
{\CC{[PSY]} \emph{Spezielle Patientengruppe:} Psychiatrie}
\label{chp:PSY}

\ChapterIntro
{%
~~~
}
{%
\begin{description}
  \item[Maintainer:] Sebastian Gabriel
  \item[Autoren:] Diverse
  \item[Reviewer:] Standard-Reviewprozess
  \item[Version:] {\DocVersion} ({\DocVersionInfo})
  \item[SHA1:] {\DocGitCommit}
\end{description}

{\TxTeilkapitelAASS}
}

\nocite{WunnPsych1,OHCSP7}
\label{topic:unterbringungsgesetz}
\dictum[Otto Pötzl, nach \cite{Jolesch2}]{Der, der den
Schlüssel hat, ist der Arzt.}

\Layout{2}{Einleitung}
{%
Die Psychiatrie ist ein sehr kompliziertes und
anspruchsvolles Fachgebiet der Medizin. Sie beschäftigt
sich einerseits mit handfesten, körperlich bedingten
Störungen, die Auswirkungen auf das Denken und die geistige
Leistungsfähigkeit haben. Andererseits hat man es in der
Psychiatrie sehr oft mit Störungen zu tun, von deren
Ursache die Wissenschaft kaum eine Ahnung hat. In der
Psychiatrie ist es auch besonders schwierig, zwischen
"`gesund"' und "`krank"' treffsicher zu unterscheiden. 
}{

}
{./images/wikipedia-pd/van_Gogh_Cafeterasse_bei_Nacht.jpg}
{Vincent van Gogh, Caféterasse bei Nacht.}{}


Grundsätzlich muss aber gesagt werden, dass psychiatrische
Erkrankungen\footnote{Sofern der Begriff
\emph{"`Erkrankung"'} überhaupt angebracht ist \ldots}
etwas alltägliches und nichts außergewöhnliches sind. Je
nach Studie leiden bis zu 30\,\% der Bevölkerung an
psychiatrischen Erkrankungen, allen voran an Depressionen
und der Alkoholkrankheit.


\dictum[Redensart]{If there's no problem, don't fix it.}

\paragraph{Der "`Krankheitsbegriff"' in der Psychiatrie} Man
trifft oft auch nicht klar umrissene Krankheitsbilder, und
auch deren jeweiliger Krankheitswert ist oft umstritten,
sehr vom einzelnen Patienten und seiner Situation abhängig:
Während der eine Mensch an seiner Depression leidet,
verarbeitet ein anderer seinen Weltschmerz zu einem Lied und
sichert sich dank der AKM\footnote{\Z{Autoren, Komponisten
und Musikverleger. Die AKM ist die größte Urheberrechts- und
Verwertungsgesellschaft in Österreich.}} eine
Einnahmequelle. (Subjektiv betrachtet erscheint die
Depression manchmal als ein treibender Motor der
Singer-/Songwriter-Szene.)




Einer der bekanntesten psychiatrischen Patienten war der
niederländische Maler \E{Vincent van Gogh}
({\sffamily\textborn}\,1853\,--\,{\sffamily\textdied}\,1890).
Im Alter von 35 Jahren erkrankte er an einer psychiatrischen Erkrankung die
von Wahnvorstellungen, Albträumen, Depressionen und
wiederholten Selbstmordversuchen geprägt war. Es sollte
nicht bei den Versuchen bleiben, doch bis zu seiner
Selbsttötung am 29. Juli 1890 schuf er hunderte Gemälde und
über tausend Zeichnungen. Heute gilt van Gogh als
bedeutender Mitbegründer und Impulsgeber der modernen
Malerei. \cite{wikiVanGogh}

\Fazit{Die psychiatrische Diagnose verhält sich zum Patienten wie ein Barcode
zur Tafel Schokolade: 
Niemand kauft eine Packung wegen des Codes.}

\begin{figure}[bh]
\begin{center}
\includegraphics[width=0.75\linewidth]{./images/wikipedia-pd/van_Gogh_Starry_Night_Over_the_Rhone.jpg}
\Captionof{figure}{Krank? \SourcePicture{Vincent van Gogh: Sternennacht über der
Rhone (September 1888).}}
\end{center}
\end{figure}


Der Krankheitswert eines Zustandes beginnt dort, wo dieser
Zustand für den Patienten selbst oder für seine Umgebung ein
Problem darstellt. Dabei gibt es keine scharfe Grenze und
die Grauzone ist beträchtlich.

\Layout{2}{Psychiatrische Erkrankung in der Praxis}
{%
In der Gesellschaft (repräsentiert durch die Gesetzgebung)
gibt es die Übereinkunft, dass bestimmte Krankheitsbilder
bzw. Symptome als behandlungswürdig bzw. sogar
behandlungspflichtig angesehen werden. Als
behandlungspflichtig gelten zum Beispiel die \E{Eigen-} und
\E{Fremdgefährdung} aufgrund einer psychiatrischen
Erkrankung, hier hat der Gesetzgeber Maßnahmen nach dem
\T{Unterbringungsgesetz} (\T{UbG}) vorgesehen. Auch der
geistige Verfall (z.\,B. im Rahmen einer (Alters-)Demenz)
kann rechtliche Schritte nach sich ziehen, im Sinne einer
Besachwaltung.
}{←}{}{}{}



\Layout{0}{Der psychiatrische Patient ist nicht entrechtet!}
{%
Es ist hierbei jedoch festzuhalten, dass auch ein
psychiatrischer Patient -- genauso wie jeder andere Patient
auch -- über die selben Selbstbestimmungsrechte verfügt,
d.h. es ihm grundsätzlich zusteht eine Behandlung zu
verweigern. Sollten im Rahmen der gesetzlichen Bestimmungen
(\EE{Unterbringungsgesetz}) Maßnahmen durchzuführen sein, so
sind diese von der \EE{Exekutive} (Polizei) vorzunehmen. Das
Ziel des Unterbringungsgesetz ist es, das Maß der
\E{Einschränkung der Freiheit des Patienten so gering wie
möglich zu halten}.

Auch im Rahmen der
\EE{Besachwaltung} kann ein \E{Gericht} das Selbstbestimmungsrecht des
Patienten einschränken. 
}{
\begin{itemize}
  \item Der psychiatrische Patient hat grundsätzlich die gleichen Rechte!
  \item Besonderheiten:
  \begin{itemize}
    \item Unterbringungsgesetz (UbG)
    \item Besachwaltung
  \end{itemize}
\end{itemize}
}{}{}{}





\Fazit{Der Rettungsdienst führt von sich aus keine
Zwangsmaßnahmen bei psychiatrischen Patienten durch.}

Für die Arbeit im Rettungsdienst ist die Kenntnis einiger
spezieller Krankheitsbilder und Symptome notwendig, sowie
ein allgemeines Verständnis über den Umgang mit
psychiatrischen Patienten.

\section{Umgang mit psychiatrischen Patienten}
\label{topic:umgang-psy}

Wie in der Einleitung erwähnt, ist eine psychiatrische
Erkrankung etwas durchaus Alltägliches, oft stellt selbige
nicht einmal den Behandlungs- bzw. Transportgrund dar.
Ebenso wurde erwähnt, dass der psychiatrische Patient
grundsätzlich über die selben (Selbstbestimmungs-)Rechte
wie jeder andere Patient verfügt. 

\paragraph{Haltung gegenüber dem Patienten} Dementsprechend
gelten grundsätzlich auch die gleichen Regeln für die
Gesprächsführung wie sie unter
\prettyref{topic:kommunikationsregeln} besprochen wurden. Der
Patient verdient den gleichen Respekt und die gleiche
Wertschätzung wie ein körperlich Erkrankter.

\paragraph{Auftreten} 

\begin{itemize}
    \item Besonders wichtig ist ein \EE{ruhiges,
    sachliches Auftreten}.
    \item  Eventuelle \EE{Aufregung} (durch den
    Patienten; Angehörige; Feuerwehr; Polizei;
    Spezialeinsatzgruppe mit Rammbock, Tränengas,
    Maschinengewehr und Angst; \ldots) \EE{soll nicht auf
    den Helfer abfärben}.
    \item Der Patient soll sich \EE{nicht beengt},
    bedroht und -- wenn möglich -- nicht bevormundet
    \EE{fühlen}.
    \item Es sollen \EE{keine falschen oder nicht
    einlösbare Versprechen} gegeben werden.
   \item \EE{Keine Lügen!}
\end{itemize}

\paragraph{Eigenschutz} Auch wenn vom durchschnittlichen
psychiatrischen Patienten keine höhere Gefahr als von der
Normalbevölkerung ausgeht, so ist es doch ratsam
grundlegende Punkte des Eigenschutzes zu beachten:

\begin{enumerate}
    \item Fluchtweg freihalten
    \item Rücken freihalten
    \item Patienten nicht beengen/bedrängen
    \item Patienten über Handlungen informieren und evtl.
    extra um Einverständnis fragen (angurten, \ldots)
\end{enumerate}


Diese Punkte sind grundsätzlich für jeden Einsatz sinnvoll!



\section{Symptome}

\subsection{Der Patient mit \emph{Wahnvorstellungen}}

Menschen können -- aus unterschiedlichen Gründen --
Wahnvorstellungen (Halluzinationen) entwickeln. Diese
können sehr unterschiedlich sein und alle Sinne betreffen.
Man kann Dinge \E{hören} (Stimmen hören, \ldots)  oder
Sachen \E{sehen} ("`weiße Mäuse"', \ldots) die nicht
existieren. Der Betroffene nimmt dies jedoch als vollkommen
"`echt"' und real war. 

Der Wirklichkeitsbezug kann teilweise erhalten oder nahezu
komplett verloren gegangen sein. 



\minisec{Fallgeschichte}

\begin{quotation}
Eine 37-jährige Patientin kommt aufgrund eines
grippalen Infektes während der Wintermonate in die
Ordination eines Arztes für Allgemeinmedizin. Während des
Gespräches erzählt die Patienten, dass sie schon seit
mehreren Jahren arbeitsunfähig sei, da sie an
Halluzinationen leide. Jedes Mal wenn sie das Badezimmer
betritt, sieht sie eine ihr unbekannte Frau beim Fenster
hinein schauen, obwohl sie im zweiten Stock wohnt. Sie
erzählt, dass sie wisse, dass es nicht real sein kann,
trotzdem nimmt sie diese Frau so wahr als würde sie
tatsächlich vor dem Fenster stehen. Sie ist nun schon seit
Jahren in psychiatrischer Behandlung, allerdings stellt
sich keine Besserung ihrer Symptome ein. Die Patientin gibt
an einen massiven Leidensdruck zu haben.
\end{quotation}



\paragraph{Fazit}

Auch psychiatrische Patienten haben "`normale"'
Alltagsprobleme, wie z.\,B. einen grippalen Infekt. Die
Wahnvorstellungen werden von den Patienten komplett real
erlebt, auch wenn sie vielleicht sogar wissen, dass sie nicht
real sein können. 

Diese Patientin hat noch eine teilweise
intakte Realitätswahrnehmung: Psychiatrische Erkrankungen
oder Symptome funktionieren \E{nicht} nach dem
"`Alles-oder-Nichts"'-Prinzip.


\paragraph{Wahnidee -- Wahnsysteme}

Auch die Wahrnehmung der Umgebung insgesamt kann gestört
sein, auch hier gibt es verschiedene Abstufungen. 

Die einfachste Form ist die \T{Wahnidee}. Hierbei hat der
Patient eine klar abgegrenzte Vorstellung von etwas, was mit
der Realität nicht vereinbar ist. 

\begin{quotation}
\emph{"`Das
Gebühren-Info-Service des ORF kontrolliert jede Woche am
Verteilerschrank, ob das Antennenkabel korrekt vom
Verteiler abgesteckt wurde."'}
\end{quotation}

Diese Wahnideen können in weiterer Folge zu einem
\T{Wahnsystem} ausgebaut werden 



\begin{quotation}
\emph{"'Das Gebühren-Info-Service
des ORF ist eine Deckmantelorganisation des israelischen
Geheimdienstes Mossad, welcher im Auftrag der CIA geheime
Abhöranlagen im ganzen Wohnhaus installiert hat. Bei
unliebsamen Äußerungen der Bewohner werden Strahlen, welche
das Denken beeinflussen können, in die Wohnungen geleitet.
Um diesen zu entgehen hat der Patient das Antennenkabel am
Verteilerkasten abgesteckt und zahlt keine Fernsehgebühren
mehr. Da ihn nun der Geheimdienst nicht mehr erreichen
kann, haben sie vom Gebühren-Info-Service einen Agenten
vorbei geschickt, welcher nun gewaltsam die Gedanken des
Patienten wegnehmen wollte."'}
\end{quotation}

\paragraph{Wahnthemen} \Z{Ein Wahn kann unterschiedliche
Themen haben, die häufigsten sind z.\,B. Verfolgungswahn,
Größenwahn oder der Verarmungswahn.}

\MassnahmenNG{Spezielle Maßnahmen}{Umgang mit Patienten mit
Wahnvorstellungen}{}{ Die -- für den Außenstehenden -- wirren Vorstellungen
stellen das Fachpersonal vor die Frage, wie mit diesen
Patienten umzugehen ist.
Ziel ist:
\begin{enumerate}
  \item Abwendung von unmittelbaren Gefahren für den
  Patienten und die Umgebung
  \item Korrekte und professionelle Durchführung eines
  Transportes in eine geeignete Einrichtung.
\end{enumerate}
Es gelten somit folgende Grundsätze:
\begin{enumerate}
  \item Sofern eine Gefährdung für den Patienten oder für
  Dritte besteht, ist diese (wenn möglich) zu beseitigen,
  dazu zählt u. U. auch der Rückzug und die Herbeiziehung
  der Exekutive.
  \item Wahnvorstellungen sollen \EE{nicht ausgeredet}
  werden.
  \item  Wahnvorstellungen sollen aber auch \EE{nicht
  bestätigt} werden.
  \item Am besten man nimmt die Schilderungen des
  Patienten \EE{interessiert, aber neutral und möglichst
  kommentarlos} zur Kenntnis (\emph{"`Mhm!"', "`Aha."'}).
\end{enumerate}
}

% \subsubsection{Realitätsverlust}
% 
% \paragraph{Mögliche Ursachen}
% 
% Endogene Psychosen (Schizophrenie,  schizoaffektive Psychose)
% 
% Organische Psychosen (Alkoholparanoia,  Drogenpsychose (z.\,B. Coke bugs
% ),  Encephalitis)
% 
% Realitätsverlust - Wahn
% 
% Wahn = falsche, allerdings in sich  logische Auffassung der Realität,
% an der  unkorrigierbar festgehalten wird  (paranoid = wahnhaft) 
% 
% Themen: Verfolgungswahn, Größenwahn,  Verarmungswahn,...


% \subsubsection{Realitätsverlust - Halluzination}
% 
% 
% Halluzinationen: Trugwahrnehmungen,  für alle Sinne möglich
% (weiße Mäuse,  Stimmen,
% "`Geisterberührungen"',...)
% 
% 
% Psychomotorische Unruhe /  Aggressivität
% 
% 
% Organische Ursachen ausschließen!!!
% 
% {\mdseries\color[rgb]{0.0,0.0,0.5019608}
% Intoxikation durch Alkohol, Drogen,  Medikamente}
% 
% 
% Hypoglykämie
% 
% St.p. SHT
% 
% Hirntumor
% 
% psychomotorische Unruhe /  Aggressivität
% 
% Schwere Erregungszustände:
% 
% Katatonie : Starre mit Sprechstereotypien, Grimassieren,
% Befehlsautomatismen, Patient imitiert Untersucher
% 
% Manie: Distanzlosigkeit, aggressiv gefärbte Agitiertheit
% 
% psychoreaktiv: abnorme Persönlichkeit
% 
% psychosoziale Krise: Reaktion auf Anlaß / Streßreaktion
% 
% psychomotorische Unruhe /  Aggressivität
% 
% Hysterie (funktionelle Erregungszustände):
% 
% anlaßgebunden, entsprechende Grundpersönlichkeit
% 
% psychosomatischer Funktionsverlust (Lähmungen, Erblindung) ohne echte
% organische Ursache
% 
% Entspricht psychischer Abwehrreaktion
% 
% \subsubsection{hysterischer Krampfanfall:}
% 
% {\mdseries
% keine Verletzungen, keine Apnoe, atypische Krämpfe, kein
% Stuhl-/Harnverlust}
% 
% {\mdseries
% unbewußt "`gespielt"'}



\subsection{\Z{Reserviert}\A{Der \emph{stuporöse}
Patient}}

%\Definition{Stupor: Erschlaffung}
% 
% \paragraph{\TxSymptome}
% 
% \begin{itemize}
%     \item ausdrucksloser Patient
%     \item unbeweglich
%     \item starre Mimik
%     \item spricht nicht, ißt nicht ={\textgreater} ev. bis Verhungern!
%     \item Suchterkrankungen
%     \item primär psychisches Problem
%     \item sekundär körperliche und soziale Folgen
%     \item eigengesetzlicher Ablauf
%     \item zunehmender Kontrollverlust über  eigenes Verhalten
% \end{itemize}

% \subsubsection{Manie}
% 
% {\mdseries
% rastlose Unruhe - "`energiegeladen"'}
% 
% {\mdseries
% gereiztes Verhalten}
% 
% {\mdseries
% Überaktivität}
% 
% {\mdseries
% Distanzverlust}
% 
% {\mdseries
% Verwirrungszustände}
% 
% {\mdseries
% Aggression}
% 
% {\mdseries
% Selbstüberschätzung}
% 
% {\mdseries
% ev. als bipolare Störung abwechselnd mit  Depression}


\subsection{Psychomotorische Unruhe und Aggressivität}

Unruhe und Aggressivität kann ein Anzeichen für viele
Krankheitsbilder sein. Die Differentialdiagnosen umfassen
z.\,B. Blutzuckerstörungen, Demenz, Intoxikation,
Panikattacken, Manie, akute Psychosen u.\,v.\,a.\,m. Hier soll nur
auf die grundsätzlich zu bedenkenden Maßnahmen im Umgang mit diesen Patienten
eingegangen werden. 

\MassnahmenNG{Spezielle Maßnahmen}{Unruhiger oder aggressiver Patient}{}{

\begin{itemize}
  \item Ruhe vermitteln und deeskalieren:
  \begin{itemize}
  \item  Ruhig bleiben! Achtung auf eigenen Tonfall!
  \item Nicht provozieren!
  \item Aufregung nicht abfärben lassen
  \item \emph{"`Talk down"'}
  \end{itemize}
  \item Selbstschutz, Fluchtweg offen halten
\end{itemize}

}

\subsection{Suizidalität}
\label{topic:suizidalität}
\label{topic:suizid1}

\Definition{\T{Suizid}: absichtliche
Selbsttötung. \T{Erweiterter Suizid}: Tötung der eigenen und
fremder Personen.}

Selbstmorde und Selbstmordversuche (\T{SMV}) gehören zu den
häufigen Einsatzgründen im Rettungsdienst. Häufig sind
Suizid- oder sonstige selbstschädliche Handlungen nicht auf
den ersten Blick als solche erkennbar. 

Wenn der Verdacht auf eine suizidale Handlung besteht, soll
dies vertraulich, aber \EE{direkt angesprochen und geklärt
werden}. Es ist dabei wichtig, einen \Ca{neutralen Tonfall}
zu verwenden und \Ca{nicht anklagend zu klingen!}

\begin{quotation}
\Speech{"`Wollten Sie sich verletzen?"'}

\Speech{"`Wollten Sie sich umbringen?"'}
\end{quotation}

Häufig berichten Patienten über Gefühle der
Hoffnungslosigkeit, Einsamkeit, Verzweiflung, innerer Leere
oder lassen einen fehlenden Realitätsbezug erkennen.
\EE{Die geäußerten Aussagen und Gefühle des Patienten sollen nicht
gewertet, verharmlost oder verniedlicht werden!}

\Cave{Jede Suizidäußerung des Patienten muss ernst
genommen und dokumentiert werden.}

Ein misslungener -- oft zögerlicher -- Selbstmordversuch
wird oft von Unwissenden im Sinne von \emph{"`Der hat es gar nicht so gemeint"'}
aufgefasst. Das ist falsch\footnote{Zumindest meistens.}!
Ein Selbstmordversuch ist oft der letzte Hilferuf, den ein
Mensch -- sei es aus Vereinsamung oder aus anderen Gründen
-- aussenden kann. 

\Cave{Jeder Selbstmordversuch muss ernst genommen
werden, auch wenn er noch so lächerlich oder unglaubwürdig
aussieht.}


\subsection{Der \emph{unmittelbar eigen- oder
fremdgefährdende} Patient}
\label{topic:unterbringung}
\label{topic:suizid2}

\Layout{22}{Zwei Möglichkeiten:}
{
\begin{enumerate}
  \item Freiwillige Behandlung (Einwilligung des Patienten)
  \item Unterbringung (Behandlung gegen oder ohne Willen
  des Patienten) gem. Unterbringungsgesetz
\end{enumerate}


Die \index{Unterbringung}Unterbringung wird durch das
\index{Unterbringungsgesetz}Unterbringungsgesetz geregelt.
Auszug:

{\S}\,3. In einer Anstalt darf nur untergebracht werden, wer
\begin{itemize}
  \item 1. an einer \EE{psychischen Krankheit} leidet und im Zusammenhang
  damit \EE{sein Leben} oder seine \EE{Gesundheit} oder das Leben
  \EE{oder} die Gesundheit \EE{anderer} ernstlich und erheblich
  gefährdet und
  \item 2. \EE{nicht in anderer Weise}, insbesondere außerhalb einer
  Anstalt, ausreichend ärztlich behandelt oder betreut werden kann.
\end{itemize}

}{}
{./images/shared/mann_mit_pistole.jpg}
{Sie kommen in eine Wohnung und finden diese
Situation vor \ldots Was machen Sie?}
{Gabriel, AASS-CC-AT.at}





\paragraph{Unterbringung ohne Verlangen}~

{\S}\,8. Eine Person darf gegen oder ohne ihren Willen nur
dann in eine Anstalt gebracht werden, wenn ein im
\EE{öffentlichen Sanitätsdienst stehender Arzt} oder ein
\EE{Polizeiarzt} sie untersucht und bescheinigt, dass die
Voraussetzungen der Unterbringung vorliegen. In der
Bescheinigung sind im einzelnen die Gründe anzuführen, aus
denen der Arzt die Voraussetzungen der Unterbringung für
gegeben erachtet.

{\S}\,9\,(1) Die \EE{Organe des öffentlichen
Sicherheitsdienstes} sind berechtigt und verpflichtet, eine
Person, bei der sie aus besonderen Gründen die
Voraussetzungen der Unterbringung für gegeben erachten, zur
Untersuchung zum Arzt ({\S}\,8) zu bringen oder diesen
beizuziehen. Bescheinigt der Arzt das Vorliegen der
Voraussetzungen der Unterbringung, so haben die Organe des
öffentlichen Sicherheitsdienstes die betroffene Person in
eine Anstalt zu bringen oder dies zu veranlassen. Wird eine
solche Bescheinigung nicht ausgestellt, so darf die
betroffene Person nicht länger angehalten werden.

(2) \EE{Bei Gefahr im Verzug} können die \EE{Organe des
öffentlichen Sicherheitsdienstes} die betroffene Person
\EE{auch ohne Untersuchung und Bescheinigung in eine Anstalt
bringen}.

(3) Der Arzt und die Organe des öffentlichen
Sicherheitsdienstes haben unter möglichster Schonung der
betroffenen Person vorzugehen und die notwendigen
Vorkehrungen zur Abwehr von Gefahren zu treffen. Sie haben,
soweit das möglich ist, mit psychiatrischen Einrichtungen
außerhalb einer Anstalt zusammenzuarbeiten und
\EE{erforderlichenfalls den }\EE{örtlichen Rettungsdienst
beizuziehen}.

Daraus ergibt sich bei \EE{fehlender}\EE{Freiwilligkeit}:

\paragraph{Voraussetzungen für eine Unterbringung ({\S}\,3)}

\begin{itemize}
    \item \EE{Psychiatrische Erkrankung}
    (ursächlich!), (${\rightarrow}$~{\S}\,3~Abs.\,1)
    \item \EE{Eigen-}  und/oder
    \EE{Fremdgefährdung},
    (${\rightarrow}$~{\S}\,3~Abs.\,1)
    \item \EE{Keine weniger einschränkende
    Möglichkeit} zur Abwendung der Gefährdung,
    (${\rightarrow}$~{\S}\,3~Abs.\,2)
\end{itemize}

\paragraph{Vorgehensweise ({\S}{\S}\,8, 9)}

Ein "`im öffentlichen Sanitätsdienst stehender Arzt
oder ein Polizeiarzt"' die Voraussetzungen zu
bescheinigen und zu Begründen. In Wien ist das der
"`\index{Amtsarzt}\EE{Amtsarzt}"',
welcher über die Polizei verständigt wird. 

\EE{Der Notarzt, auch wenn er von der Gemeinde Wien
gestellt wird, hat nicht die Kompetenz eine Unterbringung anzuordnen}.
(${\rightarrow}$~{\S}\,8)

Für den Transport ist grundsätzlich die Polizei zuständig
(${\rightarrow}$~{\S}\,9~Abs.\,1)! Die Polizei darf des Rettungsdienst
quasi als "`Taxi"' beiziehen
(${\rightarrow}$~{\S}\,9~Abs.\,3). 

Bei Gefahr im Verzug darf die Polizei auch ohne Bescheinigung eine
Person auf eine psychiatrische Abteilung zwecks Unterbringung
verbringen (${\rightarrow}$~Abs.\,2). 

Im Zweifelsfall Polizei nachalarmieren...

\paragraph{Weiters: }~

\EE{Eine Unterbringung kann nur in einer psychiatrischen Abteilung erfolgen!}

Eine Unterbringung \EE{kann auch freiwillig}
erfolgen. Da dann auch der Transport freiwillig
stattfindet, sind dabei keine besonderen Bestimmungen zu
beachten.







\section{Krankheitsbilder}

\subsection{Psychose}
\index{Psychose}
\label{topic:psychose}
\TODO{EMHO}{2010-08-04}{Text fehlt}{}

\TakeNotesLarge

\subsection{Demenz}

\Definition{Fortschreitende Einschränkung der  geistigen
Leistungsfähigkeit des Gehirns. Die häufigste Form ist die
Altersdemenz.}

\Layout{2}{Betrifft}{%
\begin{itemize}
    \item Gedächtnis
    \item Denkvermögen
    \item Sprache
    \item Motorik
    \item Persönlichkeitsstruktur
\end{itemize}
}{}{}{}{}

\Layout{2}{\TxSymptome}{%
\begin{itemize}
    \item Zeitliche und örtliche Desorientierung
    \item psychomotorische Unruhe
    \item Personenverkennung
    \item Wechselnde Symptomatik!
\end{itemize}
}{}{}{}{}



\subsection{Depression und Manie}
\label{topic:manie}
\label{topic:depression}

\TODO{GABS}{2010-08-04}{fehlt: Depression}{}

\TakeNotesLarge

 

% TODO: Depressionen 
% \subsubsection{\Z{Reserviert}\A{Depression}}
% 
% {\mdseries
% Traurigkeit}
% 
% {\mdseries
% Schuldgefühle}
% 
% {\mdseries
% fehlender Antrieb}
% 
% {\mdseries
% Schlafstörungen}
% 
% {\mdseries
% Suizidgedanken}
% 
% \paragraph{Depression / Verzweiflung / Suizidalität}~
% 
% 
% {\mdseries
% Grundstimmung entweder:}
% 
% {\mdseries
% antriebslos  oder}
% 
% {\mdseries
% agitiert (höhere Suizid gefahr!)}
% 
% \paragraph{Maßnahmen:}
% 
% {\mdseries
% auf Gegenwart konzentrieren, soziale  Umgebung miteinbeziehen, 
% Einfühlungsvermögen!!!}
% 
% {\mdseries
% Depression / Verzweiflung / Suizidalität}
% 
% {\mdseries
% Selbstmordversuch ={\textgreater} stationäre  Aufnahme, falls
% notwendig Parere!!!}
% 
% {\mdseries
% Selbstmordphantasien: je konkreter,  desto gefährlicher!}
% 
% {\mdseries
% Suizid = Selbstmord}
% 
% {\mdseries
% Achtung, evtl. sogar Mitnahmeselbstmord  möglich!}
% 
% {\mdseries
% Angst}
% 
% {\mdseries
% Gefühl der Bedrohung mit vegetativer Begleitsymptomatik}
% 
% {\mdseries
% oft ohne erkennbaren Grund}
% 
% {\mdseries
% bei psychiatrischen und organischen  Krankheitsbildern}
% 
% {\mdseries
% Angst}
% 
% 
% 
% 



\subsection{Alkohol- und Drogenentzug}
\label{topic:alkohol-entzug}


\Layout{2}{\TxSymptome}
{%
\begin{itemize}
  \item Schwitzen
  \item Unruhe
  \item Zittern (Tremor)
  \item Delirium  (Bewusstseinsstörungen, Halluzinationen )
  \item Krampfanfälle
  \item RR-Entgleisung
  \item \E{"`Cold turkey"'} bei Heroinentzug
\end{itemize}
}{

}{}{}{}

\Layout{2}{Anamnese}
{
\begin{itemize}
  \item Plötzlich kein Zugang mehr zu Alkohol:
  \item Kürzlicher Einzug in Altersheim
  \item kürzlich pflegebedürftig geworden
  \item \ldots
\end{itemize}
}{}{}{}{}












%%%%%%%%%%%%%%%%%%%%%%%%%%%%%%%%%%%%%%%%%%%%%%%%%%%%%%%%%%%

\paragraph{Entzugsdelir} Im Extremfall leidet der Patient an
einem Entzugsdelir (Delirium tremens). Dabei kann  es zu
Zittern, Unruhe, Angst, \EE{Halluzinationen} ("`weiße
Mäuse"', \ldots), vegetativer Entgleisung (Tachykardie,
Hypo-/Hypertonie) und zerebralen Krampfanfällen kommen. Es
handelt sich um einen \Ca{lebensbedrohlichen Zustand}!



\subsection{Panikattacke}

Eine Panikattacke kann sich sehr unterschiedlich äußern und
körperliche Beschwerden vortäuschen.

\begin{itemize}
    \item Plötzliches Angstgefühl
    \item Dauer: Sekunden bis Minuten
    \item Thoraxbeschwerden (Infarktangst)
    \item Hyperventilation \index{Hyperventilation!bei Panikattacke}
    \item Mehrfach durchuntersuchte  Patienten, wirken gesund
    \item Erwartungsangst
    \item Vermeidungsverhalten
    \item Antriebslosigkeit / Stupor
\end{itemize}



% TODO Schizophrenie
% \subsubsection{\Z{Reserviert}\A{Schizophrenie}}














% 
%  {\mdseries 10\% aller psychiatrischen Erkrankungen 
% liegt eine organische Ursache zugrunde.}
% 
% {\mdseries
% Bei Erstmanifestation einer  psychiatrischen Erkrankung MUSS  daher eine
% organische Ursache  ausgeschlossen werden!!!}
% 
% {\mdseries
% Der Umgang mit psychiatrischen Patienten ist heikel und erfordert viel
% Feingefühl sowie die Fähigkeit, auf das Gegenüber einzugehen. }




% 
% \subparagraph*{
% Situation:}
% 
% {\mdseries
% Belassung? (ist Patient überhaupt  reversfähig?)}
% 
% {\mdseries
% Unterbringung? (notwendig / rechtlich  zulässig?)}
% 
% {\mdseries
% Polizei / Amtsarzt bzw. Notarzt? (stabil?  Eskalation?)}
% 
% \subparagraph*{Leitsymptome}
% 
% \begin{itemize}
% \item Verwirrtheit, Bewußtseinseintrübung
% \item Realitätsverlust (Wahn, Halluzinationen )
% \item psychomotorische Unruhe, Aggressivität
% \item Depression , Verzweiflung, Suizidalität
% \item Angst
% \item Antriebslosigkeit, Stupor
% \item Verwirrtheit / Bewußtseinstrübung
% \item einfache Verwirrung
% \item Altersdemenz
% \item delirante Verwirrung
% \item Einfache Verwirrung
% 
% \end{itemize}
% {\mdseries
% Kann bei allen "`normalen"',
% "`körperlichen"' Notfällen, die mit
% Bewusstseins\-ein\-schränkungen ein\-her\-gehen, auftreten!}
% 
% {\mdseries
% Bsp.: Epilepsie , Commotio cerebri, Exsikkose etc.}
% 
% 
% delirante Verwirrung 
% organische Symptome:
% 
% Tachykardie
% 
% Mydriasis (Pupillen weit)
% 
% trockene rote Haut
% 
% Krämpfe
% 
% 
% z.\,B. Entzugsdelir bei Alkohol, Drogen
% 
% \subparagraph*{Maßnahmen}
% 
% 
% NAW, ad Interne Überwachungsstation (RR  instabil!)
% 
% 
% RR-Kontrollen

\section*{\protect\dsliterary{} Weiterführende Literatur}
\begin{WidePar}
\begin{multicols}{2}
\begin{labeling}{mmm}
  \item[\cite{ScholzNotfallmedizin2}] \emph{\bibentry{ScholzNotfallmedizin2}}
  \item[\cite{WunnPsych1}] \emph{\bibentry{WunnPsych1}}
  \item[\cite{BattegayHandwPsych2}] \emph{\bibentry{BattegayHandwPsych2}}
%   \item[\cite{DahmerGespraechsfuehrung5}] \emph{\bibentry{DahmerGespraechsfuehrung5}}
  \item[\cite{ThunRedenI46}] \emph{\bibentry{ThunRedenI46}}
  \item[\cite{ThunRedenII30}] \emph{\bibentry{ThunRedenII30}}
  \item[\cite{ThunRedenIII19}] \emph{\bibentry{ThunRedenIII19}}
  \item[\cite{ArgelanderErstinterview8}] \emph{\bibentry{ArgelanderErstinterview8}}
%   \item[\cite{}] \emph{\bibentry{}}
\end{labeling}
\end{multicols}
\end{WidePar}
